\scanpage{95}\mspage{86}
Passons aux questions de \unsure{supports|support}. Dans le cas général, on ne peut se contenter des ombres ou supports \unsure{internes|induits} sur $L$, faute de Mag L 2. \deleted{et même des} D'abord, étudions les figures (finies).
\[
\begin{aligned}
 \deleted{X \mathrel{\vcenter{\hbox{\tikz[x=.8em,y=.5em,line join=round]{\draw (0,1)--(1,.5)--(0,0) (1,1)--(0,.5)--(1,0);}}}} Y}&\left\{\begin{aligned}
 &\deleted{\text{si }X=P_x,Y=P_y:\quad x\unsure{=|\leq|\neq}y\quad\text{\unsure{ou}}\ x\,|\circ|\,y
             \quad\unsure{<x}\quad\text{i.e. }\{x,y\}\in\unsure{\mathrm{Drap}|\mathrm{Drp}}(L)}\\
 &\deleted{\text{si }X=P_x,\ Y=S_{a,b}:\quad\text{si }x\ll S_{ab}\ \text{i.e. }x\notin[a,b]:
                x\parallel S_{ab}}\\
 &\deleted{x\in\{a,b\}\quad\text{\uline{ou}}\quad x\parallel}
 \end{aligned}\right.
\end{aligned}
\]

\deleted{\unsure{R}}Examinons la condition de compatibilité $X \mathrel{\vcenter{\hbox{\tikz[x=.8em,y=.5em,line join=round]{\draw (0,1)--(1,.5)--(0,0) (1,1)--(0,.5)--(1,0);}}}} Y$. Si $X$ disons \unsure{(ou $Y$)|($Y$)} est de la forme $P_x$, donc un lieu
\[
 \begin{aligned}
 P_x \mathrel{\vcenter{\hbox{\tikz[x=.8em,y=.5em,line join=round]{\draw (0,1)--(1,.5)--(0,0) (1,1)--(0,.5)--(1,0);}}}} \deleted{\unsure{X}}Y_{\unsure{a|\epsilon|{}}}
   &\Longleftrightarrow P_x=Y\ \text{ou}\ 
      \underbrace{P_x\,|\circ|\,Y}_{\deleted{\text{donc }\{x,y\}\in\unsure{\mathrm{Dr}}}}\\
   &\hspace{1em}\unsure{:}
 \end{aligned}
\]
\margin{(vrai pour tout lieu dans un magma}

\deleted{donc}
\[
\begin{aligned}
 P_x \mathrel{\vcenter{\hbox{\tikz[x=.8em,y=.5em,line join=round]{\draw (0,1)--(1,.5)--(0,0) (1,1)--(0,.5)--(1,0);}}}} P_y&\Longrightarrow\{x,y\}\in \mathrm{Drap}(L)\\
 P_x \mathrel{\vcenter{\hbox{\tikz[x=.8em,y=.5em,line join=round]{\draw (0,1)--(1,.5)--(0,0) (1,1)--(0,.5)--(1,0);}}}} S_{a,b}&\Longleftrightarrow P_x\,|\bullet|\,S_{ab}
              \Longleftrightarrow x\leq a\ \text{ou}\ \unsure{x|a}\geq b
\end{aligned}
\]

D'autre part, $S_{ab} \mathrel{\vcenter{\hbox{\tikz[x=.8em,y=.5em,line join=round]{\draw (0,1)--(1,.5)--(0,0) (1,1)--(0,.5)--(1,0);}}}} S_{a',b'}$ signifie que, ou bien $S_{ab}=S_{a',b'}$, ou bien $S_{ab}\parallel S_{a',b'}$, ou \deleted{$\unsure{\widetilde S|S}$} $\widetilde S_{a,b}\cap\widetilde S_{a',b'}$ contenu dans $\partial S_{a,b}\cap\partial S_{a',b'}^{\unsure{+|*}}$ et non vide. \inserted{\uline{et} $S_{ab}\,|\circ|\,S_{a',b'}$}

Dans ce dernier cas, je dis que l'intersection \inserted{$\partial S\cap\partial S'$} est réduite à un \inserted{seul} \deleted{point} élément -- s'il y en avait deux, on aurait en effet $\partial S=\partial S'$, d'où $S=S'$, cas que nous avions exclu. On trouve ainsi
\[
 \text{\unsure{la}}\quad S_{a,b} \mathrel{\vcenter{\hbox{\tikz[x=.8em,y=.5em,line join=round]{\draw (0,1)--(1,.5)--(0,0) (1,1)--(0,.5)--(1,0);}}}} S_{a',b'}\Longleftrightarrow
 \left\{\begin{aligned}
 &\text{ou}^{\text{(i)}}\ S_{a,b}=S_{a',b'}\quad\text{i.e }a=a',b=b'\\
 &\text{ou}^{\text{(ii)}}\ S_{a,b}\,|\unsure{\bullet|\circ}|\,S_{a',b'},\
       \deleted{\text{\unsure{ou}}}\ \text{i.e }a<b<a'<b'\ \text{ou}\ a'<b'<a<b\\
 &\text{ou}^{\text{(iii)}}\ a<b\deleted{\unsure{<}}\unsure{\leq|=}a'<b'
                   \quad\text{ou}\ a'<b'\leq a<b\\
 &\quad\deleted{\text{\unsure{ou bien}}}\\
 &\quad\deleted{\text{i.e. }S_{ab}\,|\circ|\,S_{a',b'}}\\
 &\quad\deleted{\text{(mais alors encore } S_{a,b}\,|\circ|\,S_{a',b'}}
 \end{aligned}\right.
\]
au total
\[
\begin{aligned}
 S_{a,b} \mathrel{\vcenter{\hbox{\tikz[x=.8em,y=.5em,line join=round]{\draw (0,1)--(1,.5)--(0,0) (1,1)--(0,.5)--(1,0);}}}} S_{a',b'}&\Longleftrightarrow S_{a,b}=S_{a',b'}\ \text{ou}\ S_{a,b}\,|\circ|\,S_{a',b'}\\
 &\qquad\deleted{\text{i.e. } a<}\\
 &\qquad\text{ce dernier cas s'écrivant}\\
 &\qquad a<b\leq a'<b'\ \text{ou}\ a'<b'\leq a<b
\end{aligned}
\]
% archive sheet number: 95

